\section{The Latin rule and the Eastern boundary}

No universal sentence about ``the Catholic rite of exorcism'' is safe unless
it names a Church, book, minister, and juridical scope. This study establishes
one necessary boundary: the 1983 \emph{Codex Iuris Canonici} governs the Latin
Church, while CCEO canon 1 says that the Eastern Code concerns all and only the
Eastern Catholic Churches unless it expressly provides otherwise about
relations with the Latin Church. Canon 2 requires the Eastern canons to be
assessed principally from the ancient law they receive or adapt. Latin canon
1172 therefore may not be silently transplanted as though it were an Eastern
canon.

\subsection{The Latin reservation}

In the Latin Church, canon 1172 makes particular and express permission of the
local ordinary necessary for licit exorcism over the possessed. Permission may
be given only to a priest marked by piety, knowledge, prudence, and integrity.
Canon 134 supplies the controlling definition. Its paragraphs 1--2 include
vicars general and episcopal vicars among local ordinaries; the exclusion in
paragraph 3 instead governs canons that attribute executive power by name to
the diocesan bishop, unless a vicar has a special mandate. Canon 1172 says
\emph{local ordinary}, not \emph{diocesan bishop}. The Catechism's shorter
episcopal summary and the 2000 instruction's rule of dependence on the
diocesan bishop do not rewrite that canonical term.
The revised \emph{De Exorcismis et Supplicationibus Quibusdam} supplies the
liturgical discipline. The 2021 amendment to canon 230 \S1 concerns stable
institution in the ministries of lector and acolyte; it neither creates an
exorcist ministry nor supplies canon 1172 permission. Since 5 June 2022,
\emph{Praedicate Evangelium}
articles 88--90 assign the Dicastery for Divine Worship and the Discipline of
the Sacraments the Apostolic See's work of regulating the liturgy, preparing
or revising typical books, confirming conference-approved translations,
recognizing legitimate adaptations, and overseeing sacramental and
sacramental discipline. This is institutional competence over the book and
its discipline. It is not the local ordinary's particular and express
permission under canon 1172, the episcopal conference's approval of a
translation or adaptation, or the institution, appointment, or authorization
of an exorcist.

The universal law identifies neither a permanent
diocesan office nor a universal intake system; concrete diocesan arrangements
therefore require their own territorial and local controls.

\begin{threestable}{Question}{What the bound source establishes}{What it does
not establish}
\textbf{Minister} & Canon 1172 restricts the permission to a priest with the
stated qualities. & Priesthood alone does not confer permission for the
reserved act.\\
\textbf{Permission} & It must be particular and express and must come from the
local ordinary. & It is not supplied by a family, prayer group, lay
association, or the priest's private judgment.\\
\textbf{Office} & A diocese may make practical arrangements for receiving and
assigning cases. & Canon 1172 does not itself create a permanent diocesan
office, title, intake system, or territorial protocol.\\
\textbf{Case judgment} & Permission regulates who may act licitly. & The grant
does not by itself prove possession or displace continuing discernment and
care.\\
\textbf{Ritual governance} & Approved liturgical law governs the ecclesial act;
the Holy See's typical-book competence, a conference's territorial approval,
and a local canon 1172 permission must be distinguished. &
Historical possession of an older book or a privately preferred formula does
not establish present authorization.\\
\end{threestable}

This separation also disciplines the word ``jurisdiction.'' Canons 1172 and
134 establish the generic category of competent authority; they do not decide
the competence of a particular officeholder in a concrete matter. Nor does
calling someone a local ordinary answer questions about territorial or
personal competence, the scope or duration of a grant, its revocation, or
applicable particular law. Those questions require the competent authority's
actual office and act, the facts of place and person, and the universal and
particular law then in force. This study supplies no case-specific legal
opinion.

The permission is not a declaration that possession has been proved. It is
particular and express authorization for a reserved ecclesial act; ordination
alone does not supply it, a lay association cannot delegate it, and a family
cannot create it by consent. The 1985 CDF letter also governs assemblies in
which no exorcism strictly so called is performed but demons are directly
addressed and their identities are sought. The 2000 instruction on healing
prayers keeps exorcism distinct from healing services and under the diocesan
bishop. These interventions close attempts to evade the reservation by
changing an event's label; they do not convert every prayer for deliverance
into the major exorcism governed by canon 1172.

Territorial implementation remains subordinate. The USCCB's public account
identifies the confirmed English book and a responsible American pattern:
medical and psychological assessment, moral certitude, consent where possible,
privacy, limited participation, and fidelity to the rite. It is official
conference guidance and an identified U.S. implementation, not universal law.
A diocese elsewhere may use a different approved translation and protocol
while remaining bound by canon 1172 and the typical book.

\subsection{What this Latin rule does not decide}

The Eastern boundary is not simply the absence of Latin canon 1172. It is a
positive legal architecture with its own scope, interpretive direction,
liturgical sources, and forms of particular law. The controlling body of
common law is the \emph{Codex Canonum Ecclesiarum Orientalium}, promulgated by
John Paul II in \emph{Sacri canones} on 18 October 1990 and binding from
1 October 1991. The jurisdiction named here is the Eastern Catholic Churches,
not the Orthodox Churches and not the Latin Church. Canon 1 says that the CCEO
concerns all and only the Eastern Catholic Churches unless a provision
expressly addresses relations with the Latin Church. A Latin canon can
therefore illuminate a comparison, but it does not become Eastern common law
through proximity, shared subject matter, or an appeal to Catholic unity.

Canon 2 gives the boundary an interpretive direction. Canons that receive or
adapt the ancient law of the Eastern Churches are to be assessed principally
from that law. Canon 3 then prevents the Code from being mistaken for a
complete liturgical manual: although it often refers to liturgical books, it
ordinarily does not decide liturgical matters, and the prescriptions of those
books are to be observed unless contrary to the canons. A present Eastern
Catholic claim may therefore require three distinct controls: CCEO common law,
the particular law of the identified Church \emph{sui iuris}, and that
Church's approved liturgical book. Silence in one layer cannot be filled by
copying the corresponding Latin provision.

Two further common-law canons establish a floor without supplying a procedure.
Canon 668 \S1 defines public divine worship as worship offered in the Church's
name by persons legitimately deputed for it and through acts approved by
ecclesiastical authority. The definition excludes the inference that private
initiative alone makes an act public worship. Yet it does not name an
exorcist, reserve a prayer, identify a permission giver, or decide whether a
particular prayer is public or private. Those propositions still require the
law and liturgical source competent for them.

Canon 867 is more directly relevant. Its first paragraph defines sacramentals
as sacred signs which, in a certain imitation of the sacraments, signify
effects---especially spiritual effects---obtained through the Church's
intercession; they dispose persons to receive the principal effect of the
sacraments and sanctify the varied circumstances of life. Its second paragraph
then gives the operative routing rule: concerning sacramentals, the norms of
the particular law of one's own Church \emph{sui iuris} are to be observed.
The canon contains no exorcism-specific minister, faculty, permission,
diagnostic process, or ritual sequence. Its remittal to particular law is not
an authorization vacuum. It is a direction to identify the Church and consult
the law that governs it.

\begin{threestable}{Question}{What CCEO common law establishes}{What remains
Church-specific}
\textbf{Scope} & Canons 1--2 establish the Eastern Catholic codal boundary and
its Eastern interpretive context. & No Latin rule is transplanted merely
because it addresses the same subject.\\
\textbf{Liturgy} & Canon 3 preserves the governing place of liturgical books;
canon 668 \S1 defines public divine worship. & The approved book, text,
minister, and authorized act must be identified for the Church concerned.\\
\textbf{Sacramentals} & Canon 867 \S1 defines the category and \S2 directs the
reader to the particular law of the proper Church \emph{sui iuris}. &
Minister, reservation, permission, competent authority, territory, and
procedure cannot be supplied generically.\\
\textbf{Currentness} & \emph{Sacri canones} controls promulgation and the
1 October 1991 effective date. & Current particular law and liturgical books
must be dated and checked for the identified Church.\\
\end{threestable}

The phrase ``particular law'' also requires care. It does not name one
identical legislator or one territorial reach for every Eastern Catholic
Church. For a patriarchal Church, canon 110 \S1 assigns legislation for the
whole patriarchal Church exclusively to its synod of bishops, while canon 150
\S\S2--3 distinguishes the reach of liturgical laws, disciplinary laws, and
synodal decisions inside and outside patriarchal territory. Canon 152 applies
the common-law provisions concerning patriarchal Churches and patriarchs to
major archiepiscopal Churches and major archbishops unless the law expressly
provides otherwise or the nature of the matter indicates otherwise. For a
metropolitan Church \emph{sui iuris}, canon 167 provides for legislation by
the council of hierarchs where common law remits a matter to the Church's
particular law, with the stated Apostolic See receipt and promulgation
conditions. These canons map possible legislative organs and the reach of
their acts. They do not prove that any one of them enacted a particular
exorcism rule.

Consequently, naming ``the bishop'' without identifying the Church, office,
law, and act is not yet a canonical conclusion. Nor would the discovery of an
Eastern prayer bearing a translated title such as ``exorcism'' establish who
may use it, under what conditions, or whether the translated word covers the
same juridical category as Latin major exorcism. An office title, a liturgical
text, a minister's ordination, and an authorization are separate evidence
objects. CCEO canons on governance in general cannot be used to manufacture an
exorcism faculty that canon 867 does not state.

The \emph{Catechism of the Catholic Church} supplies a universal catechetical
control, but it must not be misidentified as the text of Eastern canon law.
Paragraph 1673 places exorcism among the Church's authoritative public acts,
distinguishes the simple form associated with Baptism from what it calls a
major exorcism, states that the latter may be performed only by a priest with
the bishop's permission, requires prudence and observance of ecclesial rules,
and insists that illness---especially psychological illness---be
distinguished and treated medically. That teaching governs the study's
doctrinal and safety boundary. It does not rewrite CCEO canon 867 into the
wording of Latin canon 1172, identify the competent Eastern authority, settle
the juridical form of permission, or establish a uniform Eastern intake and
discernment process.

Currentness was checked at two levels. The authentic promulgation layer is
\emph{Sacri canones} and the \emph{Acta Apostolicae Sedis} Latin Code. The
Holy See's official Codes of Canon Law archive, in its dated 28 July 2026
state, presents four later CCEO amendment acts: \emph{Ad tuendam fidem}
(1998), \emph{Mitis et misericors Iesus} (2015), \emph{Ab initio} (2020),
and \emph{Competentias quasdam decernere} (2022). None is presented there as
an amendment to canon 867. No official authentic interpretation of canon 867
was located in the official materials checked. These are bounded negative
results, not a claim that every act, particular norm, or liturgical directive
of every Eastern Catholic Church has been exhausted.

This study therefore states no uniform Eastern Catholic major-exorcism rule.
It has not bound the current particular law and approved exorcistic or
initiation books of each Church \emph{sui iuris}, nor the promulgation and
territorial facts needed to apply them. It also states no present Orthodox
law or practice. A responsible Church-specific account would have to name at
least the Church \emph{sui iuris}, its juridical rank and competent
legislative authority, the exact particular-law act, its promulgation and
geographic reach, the approved liturgical book and language, the kind of
prayer involved, and the date for which the claim is made.

That research limit has a practical consequence. An Eastern Catholic should
not be directed to act on a generalized summary of Latin canon 1172 or on a
privately circulated formula. The accurate route is to the competent authority
of the person's own Church, with ordinary pastoral, medical, mental-health,
emergency, and safeguarding care continuing according to their own
competences. The absence of a published rule in this study neither grants
authority nor proves that an identified Church lacks a rule. It marks the
point at which a universal comparison must stop and Church-specific evidence
must begin.

\subsection{An actor-and-act ladder}

The present discipline becomes clearer when its authorities are arranged by
the acts for which the sources make them responsible. The universal legislator
promulgates the law governing the Latin Church. Canon 1172 states the
reservation and the qualifications of the priest to whom permission may be
given. Canon 134 identifies the juridical category denoted by ``local
ordinary.'' The Apostolic See regulates the typical liturgical book and the
discipline reserved to it. An episcopal conference prepares and approves a
complete territorial translation for submission for confirmation and may
propose the adaptations contemplated by the implementing norms. The competent
local ordinary gives a particular and express permission. These are related
acts, but none substitutes for another.

This ladder prevents several common category errors. Confirmation of an
English translation does not appoint a priest. Appointment to a diocesan role
does not prove the truth of an allegation. The availability of an approved
book does not create permission to use its reserved rite. Conversely, a
particular permission does not authorize a minister to improvise another
book, bypass the discipline governing the rite, or treat authorization as a
clinical finding. Each documentary question must be answered by the source
competent for that question.

Canon 1172's two paragraphs should not be compressed into the proposition that
any priest may act whenever he personally thinks a case warrants it.
Paragraph 1 makes permission a condition of legitimate performance and
describes it as both particular and express. Paragraph 2 limits the person to
whom the local ordinary may give that permission: a priest distinguished by
piety, knowledge, prudence, and integrity of life. Priesthood is therefore a
necessary qualification in the canon, not a self-executing grant. The canon
joins personal suitability to accountable authorization rather than offering
either one as a replacement for the other.

The adjective ``particular'' also resists an unlimited inference. The checked
universal text does not by itself settle the territorial or personal facts of
a concrete matter, the scope or duration of a grant, whether it was revoked,
or what applicable particular law requires. Nor does this study possess a
specific decree from which such details could be inferred. It can describe the
universal category and insist upon an actual competent act; it cannot convert
that description into an opinion about an unidentified minister or case.

\subsection{Office, ministry, and permission}

Historical vocabulary makes this separation especially important. Earlier
Western sources can name an order or office of exorcist, and modern diocesan
usage can call a priest ``the exorcist.'' Neither use answers canon 1172's
permission question without evidence of the competent local ordinary's
particular and express permission and its demonstrated scope. The 1972 reform of former
minor orders retained lector and acolyte as instituted ministries throughout
the Latin Church and mentioned other possible ministries in connection with a
conference request to the Apostolic See. It did not establish a universal
instituted exorcist ministry. The 2021 amendment to canon 230 \S1 altered who
may be stably instituted as lector or acolyte; it did not add exorcist to that
canon or revise canon 1172.

A practical diocesan arrangement remains distinct again. A diocese may
identify a contact, receive referrals, designate advisers, or assign a priest.
Those arrangements can matter greatly for accountability, but the universal
canon does not prescribe one permanent organizational chart. A public title
therefore proves neither that every request has been accepted nor that every
act associated with the title falls within a standing permission. On the other
hand, the absence of a public office page does not prove that a diocese lacks
an authorized means of responding. Public administrative visibility,
canonical competence, and the facts of a particular authorization are three
different evidentiary questions.

The same distinction protects the afflicted person. Authorization decides who
may perform a reserved act under the Church's discipline; it neither decides a
clinical diagnosis nor displaces continuing medical, psychological, social,
safeguarding, or pastoral care. An office does not enlarge the minister's
competence into every domain that may bear upon the person's distress.

\subsection{Boundaries against relabeling}

The 1985 CDF letter supplies a precise example of discipline extending beyond
a merely verbal label. It recalls canon 1172, forbids the faithful to use the
Leonine formula against Satan and the apostate angels, and addresses assemblies
in which, although no exorcism strictly so called is performed, demons are
directly addressed and their identities are sought. The source therefore
cannot be reduced to a rule that only an event advertised as ``major
exorcism'' is accountable. At the same time, its exact scope should not be
inflated into a prohibition of every ordinary petition to God for protection
or deliverance.

The 2000 instruction on prayers for healing supplies a related but distinct
boundary. Article 8 keeps liturgical exorcism under the diocesan bishop and
forbids inserting its prayers into the celebration of Mass, the sacraments, or
the Liturgy of the Hours. The instruction distinguishes healing assemblies
from exorcism; it does not make a healing service an alternative route around
the reservation. Read together, the two acts show why ecclesial discipline
attends to what participants actually do and claim, not only to an organizer's
chosen title.

Several categories consequently must remain visible. Ordinary Christian
prayer addressed to God is not identical with direct address to a demon.
Minor exorcisms situated in Christian initiation are not declarations that a
catechumen is possessed. A celebration does not become the reserved
liturgical act merely because it is called a prayer for healing, nor may that
label be used to insert exorcism prayers or evade the reservation. A major
exorcism is not made available to an unauthorized assembly by removing its
name from a flyer. These distinctions do not publish the content
or sequence of any rite. They identify the lines that prevent devotional
language, pastoral care, and a reserved ecclesial act from being collapsed
into one undifferentiated practice.

\subsection{One territorial implementation, not a universal template}

The United States chronology illustrates the difference between universal
law and territorial implementation. The bishops approved the English
translation in November 2014, the Holy See confirmed the final text in
December 2016, and the edition took effect in the United States on 29 June
2017. Those dates establish the history of the identified U.S. English book.
They are not the promulgation date of canon 1172, the date of the Latin typical
edition, or a universal date on which every local Church adopted the same
translation.

The USCCB account also describes diocesan protocol, referral, professional
evaluation, and the relation between the authorized minister and assistants.
It is strong official evidence for the public U.S. pastoral pattern it
describes. It does not establish that every diocese uses identical internal
forms, that a conference page grants permission in a particular case, or that
its administrative description is legislation for another territory. A
responsible comparison would identify the other conference's approved book,
the relevant implementation act, and any applicable particular law rather
than treating the U.S. account as a universal handbook.

This territorial boundary is substantive, not ceremonial. Translation
governance concerns the faithful rendering of an authentic book and legitimate
adaptation for a people. Canonical authorization concerns a minister and a
reserved act. Pastoral protocol concerns how requests are received, assessed,
and cared for. Civil and professional rules govern additional responsibilities
within their own jurisdictions. A single word such as ``approval'' can obscure
all four unless the approving authority, object, territory, and date are
named.

\subsection{The Eastern limit as a positive research rule}

The absence of a verified Eastern procedure is not merely a disclaimer. It
identifies what evidence a responsible account would need: the particular
Church \emph{sui iuris}; its common and particular law; the competent hierarch;
an approved liturgical book and language; the type of prayer under discussion;
and the date and territory in which the claim is made. Without those
coordinates, a statement about ``the Eastern rite'' has no stable juridical or
liturgical object.

This means that an Eastern Catholic should not be told that the Latin
local-ordinary rule is necessarily the rule of that person's Church merely
because both Churches are Catholic. Nor can silence in this study establish
authority to use any unverified text. Referral to the competent authority of
the person's own Church is more accurate than a confident but source-less
synthesis.

The boundary also keeps historical comparison honest. Shared biblical
language, patristic inheritance, or the word ``exorcism'' may illuminate
relationships among traditions, but none alone establishes that two present
acts have the same minister, juridical reservation, liturgical form, or
pastoral process. Catholic unity does not erase the legal and liturgical
identity of the Eastern Churches, and historical resemblance does not supply
missing present law.

\subsection{A bounded Byzantine Catholic initiation comparison}

One narrow comparison can now be made without inventing a generic ``Eastern
rite.'' It concerns Christian initiation in two identified Byzantine Catholic
jurisdictions, not a stand-alone response to an allegation of possession. The
stronger evidence comes from the Metropolitan Cantor Institute of the
Byzantine (Ruthenian) Catholic Archeparchy of Pittsburgh. Its study
\emph{The Mystery of Baptism} identifies the official Slavonic order in the
\emph{Malyj Trebnyk}, printed at Rome in 1952, at pages 25--71. It then
distinguishes three English stages: a 1952 Byzantine Seminary Press booklet
marked ``for private use only,'' which nevertheless served for many years; a
provisional 1994 Parma order based on the 1952 Slavonic book and the 1730
Greek euchologion edited by Jacobus Goar; and a revised order promulgated for
the Eparchy of Passaic in 1997. The study says that the Passaic edition became
widely used in other eparchies. That last description is evidence of reported
use, not proof that one promulgation governed every Ruthenian eparchy.

The same study locates the prayers precisely within initiation. In the fuller
adult pattern, reception into the catechumenate and instruction lead to three
prayers associated with the fourth, fifth, and sixth weeks of the Great Fast.
They are situated in the narthex and precede profession of faith, baptism,
chrismation, and Eucharistic communion. When an infant is initiated, the
catechumenal ceremonies are compressed into the immediate preparation for
baptism, and the study describes two such prayers before the sponsors'
profession of faith. The historical adult sequence and its infant adaptation
therefore belong to one initiatory architecture. Their label does not turn
them into evidence for the diagnosis, authorization, or conduct of a
stand-alone rite over a person alleged to be possessed.

A second Institute page makes the source and approval boundary unusually
clear. \emph{The Exorcisms Before Baptism} says that its English texts are
translated from the Roman 1952 \emph{Malyj Trebnyk} and based on work of the
Inter-Eparchial Liturgical Commission. It also calls them unofficial and
provisional and directs readers to obtain episcopal permission for any use
beyond private self-study. The page's headings corroborate the three-week
Great-Fast sequence. This study therefore uses the page only for provenance,
status, number, and placement. It reproduces neither its formulas nor its
performative directions.

Limited Melkite evidence supports comparison at the level of structure, not
textual identity. A pastoral explanation published by the Melkite Greek
Catholic Eparchy of Newton places several prayers called exorcisms at the
opening of the catechumen rites in the narthex. It then describes renunciation
while facing west and adherence to Christ while facing east, before the main
initiation ceremony. This is useful corroboration that a Byzantine Catholic
source outside the Ruthenian jurisdiction also embeds such material in
initiation. But the six-page explanation is catechesis, not an identified
promulgated euchologion. It cannot establish the exact wording of an approved
Melkite ritual book, a common rule for every Melkite territory, or uniformity
among Byzantine Catholic Churches.

These witnesses produce a modest but important conclusion. The word
``exorcism'' spans acts whose ecclesial location matters. In the checked
Ruthenian and Melkite sources, the material described here serves the passage
from catechumenate into baptismal life. It should not be collapsed into the
Latin major-exorcism category discussed elsewhere in this study. Conversely,
the presence of baptismal exorcism language cannot answer who may undertake a
different act, under what authorization, from which approved book, or in
response to what pastoral allegation in an Eastern Catholic Church.

The limits also protect the comparison from confessional conflation. Goar's
1730 Greek edition helps explain one stated source behind Parma's 1994 work,
but a Greek or Orthodox book does not become current Catholic law merely
because a Catholic editor or commission consulted it. Nor does shared
Byzantine inheritance erase the authority of a Church \emph{sui iuris}, its
hierarchs, its particular law, or its approved editions. A fuller account
would still need the exact current ritual book and particular law of each
named Church and territory. Until those are controlled, this subsection
establishes an initiation sequence and an edition history---not a
pan-Byzantine procedure or a stand-alone faculty rule.

\subsection{Authorization does not decide the allegation}

The actor-and-act ladder ends in a distinction that should govern the whole
current-practice account. A competent authority can decide that a qualified
priest may perform an ecclesial act; that decision is not identical with a
historian's judgment about past events, a clinician's diagnosis, or a claim
that possession has been demonstrated as an observable fact. Permission
answers a juridical question. Discernment remains an accountable pastoral and
ecclesial process, and professional assessment retains its own objects.

This prevents circular reasoning in both directions. The existence of an
authorized exorcist does not prove that a diocese has found possession in every
person referred to him. A decision not to authorize the reserved rite does not
prove that the person's suffering is unreal or unworthy of care. Likewise, a
clinician's inability to explain every reported feature does not generate
canonical permission, while a clinical diagnosis does not itself determine
what ordinary pastoral support the Church may offer.

Documentary language should therefore name the proposition actually
established. An appointment establishes an appointment. A grant establishes
permission within its demonstrated scope. A conference account establishes
the public territorial guidance it states. A typical book establishes an
approved liturgical text and discipline within its authority. None should be
made to carry the unrecorded factual conclusions of a particular case.

\subsection{Accountability after a grant}

Particular and express permission is not the end of responsibility. The
USCCB's account joins authorization to prudence, professional input,
confidentiality, limited participation, and fidelity to the rite. These
features show that accountability continues after a minister is selected.
They do not create a universal administrative code, but they prevent the grant
from being imagined as an exemption from continuing discernment and care.

The same point explains why private confidence cannot replace ecclesial
oversight. A minister may possess experience, learning, or a strong personal
conviction and still remain bound by the authority, book, territorial
discipline, and safeguards applicable to the act. Charisma asserted by an
individual is not the juridical act described by canon 1172. Apparent success,
popular demand, or testimony from participants does not retrospectively
manufacture permission.

Accountability also includes restraint in public claims. Confidential cases
cannot responsibly be converted into statistics when the underlying records,
definitions, and denominators are unavailable. A diocesan title does not
establish frequency, and an edition's publication does not establish use.
Current law can be described firmly while empirical claims about prevalence,
outcomes, or ordinary diocesan practice remain unmade. This is the evidentiary
discipline required when law, pastoral care, and vulnerable persons meet.
